% 第1章 绪论

\chapter{绪论}\label{ch:introduction}

\section{研究背景与动机}\label{sec:background}

% 工业机器人编码器故障的现实背景
% - 5 种偏差来源（详见 §3.1 Tab）：
%   出厂标定残差 ~10^-3 (实证) / 温漂 ~10^-3-10^-2 / 更换部件未重标 ~10^-1
%   / 碰撞后错位 >=10^-1 / 编码器噪声 ~10^-3
% - 现有方案的不足：
%   (i) 停机重标 → 代价高、影响产线
%   (ii) 传统 FDI/FTC → 需明确故障模型 + 以"保持稳定"非"完成任务"为目标
% - 需求：在不停机的前提下在线适应未知编码器偏差

\section{问题陈述}\label{sec:problem_statement}

% 核心观察：编码器偏差使控制问题从 MDP 退化为 POMDP
% - 单一故障源 → 三重同时效应（初始偏移、控制偏差、感知偏差）→ §3.2
% - 智能体无法观测真实关节角度 → 部分可观测 → §3.3
% - 信息论上需观测增强才能恢复可辨识性 → §3.4 (Prop 1, Prop 2)
% 额外挑战：
% - 视觉 sim-to-real gap 限制纯仿真训练 → 部署需在真机继续学习
% - 真机数据稀缺 + 稀疏奖励 → SAC online 不收敛 (§6.2.5 negative result)
% - HIL 训练引入人机安全问题 → 需备份避让策略 (§5.3, §6.3)

\section{论文贡献}\label{sec:contributions}

% 三个主要贡献（对应论文标题 Hierarchical IL+RL for Fault-Tolerant Manipulation）：
%
% (1) POMDP 形式化与可观测性分析（理论贡献，§3）
%   - 编码器偏差的因果链建模（一源三效应）
%   - 严格命题：基础观测下不可辨识 (Prop 1) / 增强观测下唯一可恢复 (Prop 2)
%   - 可辨识性 ≠ 可解性的澄清 + 真机简化的范围限定
%
% (2) 分层模仿与强化学习的容错操作框架（方法+实证贡献，§4-§6）
%   - 共享方法层：25D 仿真观测 + 随机偏差训练 + DINOv3-S 共享网络架构
%   - sim 端 HIL-SERL 验证理论：H1-H3 实验闭环
%   - real 端 BC + HG-DAgger 部署：避开 SAC 在小数据真机的 actor 漂移
%   - task ⊕ backup 双策略分层：运行时距离判据硬切换
%
% (3) 真机三任务在编码器偏差注入下的端到端验证（实证贡献，§6.4）
%   - 三任务 × 4 列联合成功率矩阵（无 bias / 有 bias / +backup × 2）
%   - HIL-SERL 真机失败的方法论级 negative result
%   - V3→V3b 几何升级是真机失效模式倒逼仿真升级的案例

\section{论文组织结构}\label{sec:organization}

% 第 2 章：相关工作（FDI/FTC、POMDP、IL/RL 家族、安全 RL + Reward hacking）
% 第 3 章：编码器偏差建模与 POMDP 分析（理论核心）
% 第 4 章：任务策略——方法设计与共享网络架构
% 第 5 章：仿真——策略训练与可观测性验证（HIL-SERL + Backup SAC+DR + 实验）
% 第 6 章：真机——模仿部署与端到端验证（BC+HG-DAgger + WiLoR 备份 + 4 列矩阵）
% 第 7 章：（已合并到第 6 章；plan B 重组将原"真机部署与验证"吸入第 6 章）
% 第 8 章：总结与展望（含 future work 两步消融）
